In two-agent LLM solver-checker systems, we uncover a perturbation-sensitivity dissociation: acknowledgments---the most individually disruptive messages---can be safely pruned in aggregate. Using a reusable analysis protocol (5-class communicative function taxonomy, perturbation-based impact measurement, non-inferiority and TOST equivalence testing), we establish this across two backbone families, \numValidatedDomains{} reasoning domains, and multiple scales (7B primary, with 14B and 72B probes), confirmed by 10-seed replication ($p < 0.001$). Three lines of correlational evidence are consistent with our \emph{cascade dissolution} conjecture: self-reinforcing ack cascades, near-zero cross-model concordance, and information-sparse acknowledgment profiles. For Llama, a unidirectional error channel---acknowledgments carry errors but never corrections---explains why pruning improves accuracy, a mechanism distinct from human backchannels requiring causal verification.

Acknowledgment pruning eliminates 36--56\% of inference calls (entire LLM invocations, not token-level compression) with non-inferior accuracy, replicated on GSM8K and StrategyQA. Boundary conditions---heuristic non-transferability, domain mismatch (MATH), stochastic decoding sensitivity---motivate empirical screening before deployment. \textbf{All results are under greedy decoding in two-agent topologies.}
